% This document was generated by an LLM.

\documentclass{essay}

\usepackage{siunitx}
\DeclareSIUnit{\var}{var}

\newcommand{\ee}{\mathrm{e}}
\newcommand{\jj}{\mathrm{j}}
\renewcommand{\Re}{\operatorname{Re}}
\renewcommand{\Im}{\operatorname{Im}}

\title{AC and Transient Analysis}
\subtitle{A Formula Reference for Linear Circuits --- Inductors, Capacitors,
Impedance, Charge, Flux, Voltage \& Current}
\author{}
\date{}

\begin{document}

\maketitle

This reference collects the constitutive relations of passive elements, the
governing equations of first- and second-order transients, and the
phasor/impedance description of sinusoidal steady state. Throughout, the
passive sign convention is used: current \(i\) enters the terminal at which the
voltage \(v\) is taken to be positive. The imaginary unit is written \(\jj\) so as
not to clash with current, and \(\omega = 2\pi f\) denotes angular frequency.

\section{Constitutive Relations of Passive Elements}

The three ideal two-terminal elements are defined by how voltage and current
relate at their terminals.

\subsection{Resistor}
A resistor obeys Ohm's law instantaneously; it stores no energy.
\[ v(t) = R\,i(t), \qquad i(t) = \frac{v(t)}{R}, \qquad p(t)=v\,i = i^{2} R = \frac{v^{2}}{R}. \]

\subsection{Capacitor --- charge and current}
A capacitor stores energy in its electric field. Its defining variable is the
stored charge \(q\), proportional to the terminal voltage through the capacitance
\(C\):
\[ q(t) = C\,v(t), \qquad i(t) = \frac{\dd q}{\dd t} = C\,\frac{\dd v}{\dd t}, \qquad v(t) = \frac{1}{C}\int_{-\infty}^{t} i(\tau)\,\dd\tau . \]
With a known initial voltage \(v(t_{0})\),
\[ v(t) = v(t_{0}) + \frac{1}{C}\int_{t_{0}}^{t} i(\tau)\,\dd\tau . \]
The voltage across a capacitor cannot change instantaneously for a finite
current, so \(v_{C}\) is a continuous state variable.

\subsection{Inductor --- flux and voltage}
An inductor stores energy in its magnetic field. Its defining variable is the
flux linkage \(\lambda\), proportional to the current through the inductance \(L\):
\[ \lambda(t) = L\,i(t), \qquad v(t) = \frac{\dd \lambda}{\dd t} = L\,\frac{\dd i}{\dd t}, \qquad i(t) = \frac{1}{L}\int_{-\infty}^{t} v(\tau)\,\dd\tau . \]
With a known initial current \(i(t_{0})\),
\[ i(t) = i(t_{0}) + \frac{1}{L}\int_{t_{0}}^{t} v(\tau)\,\dd\tau . \]
The current through an inductor cannot change instantaneously for a finite
voltage, so \(i_{L}\) is a continuous state variable.

\subsection{Stored energy}
\[ W_{C}(t) = \tfrac{1}{2}\,C\,v^{2} = \frac{q^{2}}{2C}, \qquad W_{L}(t) = \tfrac{1}{2}\,L\,i^{2} = \frac{\lambda^{2}}{2L}. \]

\begin{center}
  \begin{tabular}{@{}lccc@{}}
    \hline
    & \textbf{Resistor} & \textbf{Capacitor} & \textbf{Inductor}\\
    \hline
    \(v\)--\(i\) law      & \(v=Ri\)                 & \(i=C\,\dd v/\dd t\)
    & \(v=L\,\dd i/\dd t\)\\
    State variable    & ---                    & \(v_{C}\) (charge
    \(q=Cv\))    & \(i_{L}\) (flux \(\lambda=Li\))\\
    Stored energy     & \(0\)                    & \(\tfrac{1}{2} C v^{2}\)
    & \(\tfrac{1}{2} L i^{2}\)\\
    Cannot jump       & ---                    & voltage
    & current\\
    \hline
  \end{tabular}
\end{center}

\section{First-Order Transients}

A circuit with a single energy-storage element reduces, via Th\'evenin
reduction of the resistive part, to a first-order linear ODE. Let \(x(t)\) be the
state variable (\(v_{C}\) or \(i_{L}\)).

\subsection{General first-order form}
\[ \frac{\dd x}{\dd t} + \frac{1}{\tau}\,x = \frac{1}{\tau}\,x_{\infty}, \qquad x(t) = x_{\infty} + \bigl[x(0^+) - x_{\infty}\bigr]\,\ee^{-t/\tau}. \]
Here \(x(0^+)\) is the value just after the switching instant, \(x_{\infty}\) is the
final (steady-state) value, and \(\tau\) is the time constant.

\subsection{Time constants}
\[ \tau_{RC} = R_{\mathrm{th}}\,C, \qquad \tau_{RL} = \frac{L}{R_{\mathrm{th}}} . \]
\(R_{\mathrm{th}}\) is the Th\'evenin resistance seen by the storage element with
independent sources deactivated.

\subsection{RC circuit: charging and discharging}
For a capacitor charged through \(R\) from a source \(V_{s}\) (initially uncharged):
\[ v_{C}(t) = V_{s}\!\left(1-\ee^{-t/RC}\right), \qquad i(t) = \frac{V_{s}}{R}\,\ee^{-t/RC}. \]
For a capacitor discharging from \(V_{0}\) through \(R\):
\[ v_{C}(t) = V_{0}\,\ee^{-t/RC}, \qquad i(t) = -\frac{V_{0}}{R}\,\ee^{-t/RC}. \]

\subsection{RL circuit: energizing and de-energizing}
For an inductor energized through \(R\) from \(V_{s}\) (initially zero current):
\[ i_{L}(t) = \frac{V_{s}}{R}\!\left(1-\ee^{-Rt/L}\right), \qquad v_{L}(t) = V_{s}\,\ee^{-Rt/L}. \]
For an inductor de-energizing from \(I_{0}\) through \(R\):
\[ i_{L}(t) = I_{0}\,\ee^{-Rt/L}, \qquad v_{L}(t) = -I_{0} R\,\ee^{-Rt/L}. \]

After \(5\tau\) the response is within \(0.7\%\) of its final value and is treated
as settled.

\section{Second-Order Transients (Series \& Parallel RLC)}

With both \(L\) and \(C\) present the natural response satisfies a second-order
ODE. Writing the source-free equation for the state variable \(x\):
\[ \frac{\dd^{2} x}{\dd t^{2}} + 2\alpha\,\frac{\dd x}{\dd t} + \omega_{0}^{2}\,x = 0, \qquad s^{2} + 2\alpha s + \omega_{0}^{2} = 0, \qquad s_{1,2} = -\alpha \pm \sqrt{\alpha^{2} - \omega_{0}^{2}}. \]

\subsection{Defining parameters}
\[ \omega_{0} = \frac{1}{\sqrt{LC}} \;\;(\text{undamped natural freq.}), \qquad \zeta = \frac{\alpha}{\omega_{0}} \;\;(\text{damping ratio}), \qquad Q = \frac{1}{2\zeta}. \]
\begin{center}
  \begin{tabular}{@{}lcc@{}}
    \hline
    \textbf{Quantity} & \textbf{Series RLC} & \textbf{Parallel RLC}\\
    \hline
    Neper frequency \(\alpha\) & \(\dfrac{R}{2L}\) & \(\dfrac{1}{2RC}\)\\
    Resonant freq. \(\omega_{0}\) & \(\dfrac{1}{\sqrt{LC}}\) &
    \(\dfrac{1}{\sqrt{LC}}\)\\
    Quality factor \(Q\) & \(\dfrac{1}{R}\sqrt{\dfrac{L}{C}}\) &
    \(R\sqrt{\dfrac{C}{L}}\)\\
    \hline
  \end{tabular}
\end{center}

\subsection{Damping cases}
\begin{itemize}
  \item \textbf{Overdamped} (\(\alpha > \omega_{0}\)):
    \(x(t) = A_{1}\ee^{s_{1} t} + A_{2}\ee^{s_{2} t}\).
  \item \textbf{Critically damped} (\(\alpha = \omega_{0}\)):
    \(x(t) = (A_{1} + A_{2} t)\,\ee^{-\alpha t}\).
  \item \textbf{Underdamped} (\(\alpha < \omega_{0}\)):
    \(x(t) = \ee^{-\alpha t}\bigl(A_{1}\cos\omega_{d} t + A_{2}\sin\omega_{d} t\bigr)\),
    with damped frequency \(\omega_{d} = \sqrt{\omega_{0}^{2} - \alpha^{2}}\).
\end{itemize}

\section{Sinusoidal Steady State and Phasors}

A sinusoid \(x(t) = X_{m}\cos(\omega t + \phi)\) is represented by the phasor
\(\mathbf{X} = X_{m}\,\ee^{\jj\phi} = X_{m}\angle\phi\), recovered through
\[ x(t) = \Re\!\left\{\mathbf{X}\,\ee^{\jj\omega t}\right\}. \]
Time differentiation maps to multiplication by \(\jj\omega\), and integration to
division by \(\jj\omega\):
\[ \frac{\dd}{\dd t} \;\longleftrightarrow\; \jj\omega, \qquad \int \dd t \;\longleftrightarrow\; \frac{1}{\jj\omega}. \]

\subsection{Impedance and admittance}
Impedance \(Z = V/I\) generalizes resistance to the complex frequency domain;
admittance is its reciprocal \(Y = 1/Z = G + \jj B\).
\[ Z_{R} = R, \qquad Z_{L} = \jj\omega L, \qquad Z_{C} = \frac{1}{\jj\omega C} = -\frac{\jj}{\omega C}. \]
In rectangular and polar form,
\[ Z = R + \jj X = |Z|\,\ee^{\jj\theta}, \qquad |Z| = \sqrt{R^{2} + X^{2}}, \qquad \theta = \arctan\!\frac{X}{R}, \]
where \(X\) is the reactance: \(X_{L} = \omega L > 0\) for an inductor and
\(X_{C} = -1/(\omega C) < 0\) for a capacitor.

\subsection{Series and parallel combination}
\[ Z_{\text{series}} = \sum_{k} Z_{k}, \qquad \frac{1}{Z_{\text{parallel}}} = \sum_{k} \frac{1}{Z_{k}}, \qquad Z_{1\parallel 2} = \frac{Z_{1} Z_{2}}{Z_{1} + Z_{2}}. \]

\subsection{Resonance}
Series and parallel RLC circuits resonate when the reactive parts cancel,
\(\omega_{0} L = 1/(\omega_{0} C)\):
\[ \omega_{0} = \frac{1}{\sqrt{LC}}, \qquad f_{0} = \frac{1}{2\pi\sqrt{LC}}, \qquad \text{bandwidth } \mathrm{BW} = \frac{\omega_{0}}{Q}. \]

\section{AC Power}

For voltage and current phasors \(\mathbf{V}=V_{m}\angle\theta_{v}\),
\(\mathbf{I}=I_{m}\angle\theta_{i}\), define the phase difference
\(\theta = \theta_{v} - \theta_{i}\) and RMS magnitudes
\(V_{\mathrm{rms}}=V_{m}/\sqrt{2}\), \(I_{\mathrm{rms}}=I_{m}/\sqrt{2}\).

\[ \mathbf{S} = \mathbf{V}_{\mathrm{rms}}\,\mathbf{I}_{\mathrm{rms}}^{*} = P + \jj Q, \qquad |\mathbf{S}| = V_{\mathrm{rms}} I_{\mathrm{rms}}. \]
\[ P = V_{\mathrm{rms}} I_{\mathrm{rms}}\cos\theta \;\;(\text{real, \si{\watt}}), \qquad Q = V_{\mathrm{rms}} I_{\mathrm{rms}}\sin\theta \;\;(\text{reactive, \si{\var}}). \]
The power factor is \(\mathrm{pf} = \cos\theta = P/|\mathbf{S}|\) --- lagging for
inductive loads (\(\theta>0\)) and leading for capacitive loads (\(\theta<0\)).

\end{document}
